\chapter{The Treasury Market}\label{ch:m2:the-treasury-market}

At 9:33:19 New York time on 15 October 2014 the yield of the ten-year US
Treasury note stood at 2.02\%. Six minutes later, at 9:39:39, it was 1.86\%;
at 9:44:35 it was back at 1.99\%. No economic release, no statement, no
headline explained either move. The note is the benchmark of the dollar
yield curve, and for twelve minutes it moved like a small stock. The official report that followed found that more
than half of the trading that day, in the cash market's electronic order
books and in the futures, had been done by principal trading firms, not by
the banks that had once made the market. This chapter describes how
Treasuries are sold, where they trade and who trades them, and what two of
the market's most studied days showed about its structure.

\section{Auctions and primary dealers}

The Treasury sells its securities in auctions on a published calendar. Bids
are yields, not prices; the coupon of a new note is set only afterwards.

\begin{definition}[Uniform-price auction]\label{def:m2:the-treasury-market:auction}
In a \emph{uniform-price auction}\index{uniform-price auction} (single-price
auction) noncompetitive bids, for a limited amount, are filled first at
whatever the auction's yield turns out to be; competitive bids, each a yield
and an amount, are then accepted from the lowest yield upwards until the
offering is filled. The highest accepted yield is the \emph{stop-out} (high)
yield: every winning bidder pays it, and bids exactly at it are filled pro
rata.
\end{definition}

\begin{method}[Running an auction]\label{met:m2:the-treasury-market:run}
(1) Cap each bidder's recognised bids at 35\% of the offering, lowest yields
first. (2) Subtract noncompetitive awards from the offering. (3) Sort
competitive bids by yield and accumulate; the yield at which the cumulative
amount reaches what remains is the stop-out. (4) Fill bids below it in full;
fill bids at it by the percentage that completes the offering, rounded up to
a hundredth of a percent. (5) Set the coupon to the eighth of a percent that
gives a price closest to, but not above, par at the stop-out yield.
\end{method}

\begin{definition}[Primary dealer]\label{def:m2:the-treasury-market:pd}
A \emph{primary dealer}\index{primary dealer} is a firm that the Federal
Reserve Bank of New York has designated as a counterparty for its operations,
in exchange for obligations: to bid on a pro-rata basis in every Treasury
auction at reasonably competitive prices, to make markets for the Bank on
behalf of its official account holders, and to report on its activity.
\end{definition}

The obligation to bid makes the dealers the auction's underwriters: if
investors stay away, the dealers' bids fill it, at a yield high enough to be
worth their risk. How much of an auction the dealers had to take is therefore
read, like the auction's other statistics, as a measure of demand.

\begin{definition}[Bid-to-cover ratio and auction tail]\label{def:m2:the-treasury-market:tail}
The \emph{bid-to-cover ratio}\index{bid-to-cover ratio} is the amount bid
divided by the amount offered. The \emph{auction tail}\index{auction tail}
is the stop-out yield minus the when-issued yield
(\cref{def:m2:the-treasury-market:wi}) at the bid deadline: positive, the
auction cleared cheaper than the market (it \emph{tailed}); negative, richer
(it \emph{stopped through}).
\end{definition}

\begin{example}[A synthetic ten-year auction]\label{ex:m2:the-treasury-market:auction}
USD~42 billion of ten-year notes are offered; USD~300 million of
noncompetitive bids leave USD~41.7 billion for 76.1 billion of competitive
bids. The when-issued note trades at 4.180\% at 13:00. The bids of
\cref{fig:m2:the-treasury-market:demand} reach 41.7 billion at 4.198\%: that is
the stop-out, bids there are filled at 95.02\%, the tail is 1.8 basis points
and the \omterm{def:m2:the-treasury-market:tail}{bid-to-cover ratio} 1.82. Indirect bidders (investors bidding through a
dealer) take 51.1\%, direct bidders 10.6\% and dealers 38.3\%. The coupon is
set at 4.125\%, and every winner pays 99.409 for a note that the market valued
at 4.180\% a minute before. The auction is simulated; the procedure is the
Treasury's.
\end{example}

\begin{omfigure}
\begin{tikzpicture}
\begin{axis}[width=11cm,height=5cm,
  xlabel={yield bid (\%)},ylabel={cumulative bids (USD billion)},
  tick label style={font=\small},label style={font=\small},
  xmin=4.17,xmax=4.225,ymin=0,ymax=80,grid=major,grid style={gray!25},
  xtick={4.17,4.18,4.19,4.2,4.21,4.22},
  xticklabel style={/pgf/number format/fixed,/pgf/number format/precision=2},
  table/col sep=comma]
\addplot[omDef,very thick,const plot] table[x=yield,y=cumbn]{figdata/markets-2/04-the-treasury-market/demand.csv};
\draw[omThm,thick,dashed] (axis cs:4.17,41.7) -- (axis cs:4.225,41.7);
\draw[omProp,thick] (axis cs:4.198,0) -- (axis cs:4.198,80);
\draw[gray,thick,densely dotted] (axis cs:4.180,0) -- (axis cs:4.180,80);
\node[font=\footnotesize,anchor=north east,text=omThm] at (axis cs:4.2245,41) {to be sold: 41.7};
\node[font=\footnotesize,anchor=north west,text=omProp] at (axis cs:4.1985,76) {stop-out 4.198};
\node[font=\footnotesize,anchor=north west,text=gray] at (axis cs:4.1805,76) {when-issued 4.180};
\end{axis}
\end{tikzpicture}

\omcaption{The demand curve of \cref{ex:m2:the-treasury-market:auction}: bids
accumulated from the lowest yield. The stop-out is where the curve crosses the
amount to be sold; the gap between it and the when-issued yield at the
deadline is the tail. A book whose bids sit closer to the when-issued level
would stop through. Data: the chapter's simulated bid book.}\label{fig:m2:the-treasury-market:demand}
\end{omfigure}

\begin{dated}{2026-09}{Treasury auctions and primary dealers}\label{dat:m2:the-treasury-market:auctions}
Single-price auctions for all marketable securities since November 1998.
Noncompetitive bids up to USD~10 million per bidder, filled in full. No
competitive bidder is recognised for more than 35\% of the offering. Bids
usually close at 12:00 (noncompetitive) and 13:00 (competitive) New York
time. The Federal Reserve Bank of New York lists 26 \omterm{def:m2:the-treasury-market:pd}{primary dealers}, the most
recent added on 15 January 2026.
\end{dated}

\section{When-issued, on-the-run and off-the-run}

\begin{definition}[When-issued trading]\label{def:m2:the-treasury-market:wi}
\emph{When-issued trading}\index{when-issued trading} is trading in a security
that has been announced but not yet auctioned or issued, for settlement on
its issue date. Because the coupon is not yet known, it is quoted in yield.
\end{definition}

\begin{definition}[On-the-run and off-the-run]\label{def:m2:the-treasury-market:otr}
The most recently auctioned security of each original maturity is
\emph{on-the-run}\index{on-the-run} (the benchmark); all earlier issues of that
maturity are \emph{off-the-run}\index{off-the-run}. Each new auction makes the
previous benchmark off-the-run.
\end{definition}

\omterm{def:m2:the-treasury-market:wi}{When-issued trading} lets dealers pre-sell the auction to clients and hedge
their bids: a dealer that expects to be awarded a billion short-sells a
billion when-issued and buys it back in the auction. It also makes the tail
observable, since there is a market price at the deadline to compare with the
result. Once issued, the new note becomes the market's reference: most
electronic trading, most hedging and most repo specialness
(\cref{ch:m2:repo-and-specials}) concentrate in it. It usually yields a little
less than an \omterm{def:m2:the-treasury-market:otr}{off-the-run} note of nearly the same maturity. The difference pays
for liquidity and for the cheaper financing that its specialness brings, and
trading it, long the cheap \omterm{def:m2:the-treasury-market:otr}{off-the-run} against the rich benchmark, is one of
the oldest relative-value trades in the market (One Quant Book 9).

\section{Inter-dealer and dealer-to-client}

\begin{definition}[Inter-dealer broker and dealer-to-client platform]\label{def:m2:the-treasury-market:idb}
An \emph{inter-dealer broker}\index{inter-dealer broker} operates the venue on
which dealers, and increasingly non-dealers, trade with one another,
anonymously: for benchmark Treasuries, a central limit order book; for
\omterm{def:m2:the-treasury-market:otr}{off-the-run} issues, mostly voice-assisted brokers. A \emph{dealer-to-client
platform}\index{dealer-to-client platform} lets a client request quotes from
several dealers at once and trade on the best (\cref{ch:m2:how-bonds-trade}).
\end{definition}

The two tiers price differently. In the order books, a price is firm and
anyone who hits it trades; at the client tier, a dealer quotes a client it
knows, for a size the client names, and hedges the result in the order book.
In 2014 the inter-dealer order books for benchmark securities were BrokerTec
and eSpeed. The first offers a \emph{workup}: after a trade in the book, all
participants may trade more at that price for a short window, and the report
found that most of BrokerTec's volume happened in workups.

\begin{omfigure}
\begin{tikzpicture}[font=\small,
  box/.style={draw,thick,rounded corners=2pt,align=center,minimum height=0.95cm,inner sep=4pt},
  arr/.style={-{Stealth[length=2.5mm]},thick}]
\node[box,fill=omProp!8,draw=omProp,minimum width=2.6cm] (tsy) at (0,2.4) {Treasury\\(auctions)};
\node[box,fill=omDef!8,draw=omDef,minimum width=2.8cm] (pd) at (4.6,2.4) {primary dealers\\and direct bidders};
\node[box,fill=omThm!8,draw=omThm,minimum width=3.2cm] (idb) at (4.6,0) {inter-dealer order books\\(dealers and PTFs)};
\node[box,fill=gray!8,draw=gray,minimum width=2.8cm] (fut) at (0,0) {futures exchange};
\node[box,fill=omDef!8,draw=omDef,minimum width=2.8cm] (cli) at (9.6,2.4) {clients: funds,\\central banks, banks};
\node[box,fill=omProp!8,draw=omProp,minimum width=2.8cm] (ccp) at (9.6,0) {clearing house};
\draw[arr] (tsy) -- node[above=11pt,font=\footnotesize] {new issues} (pd);
\draw[arr,<->] (pd) -- node[right,font=\footnotesize] {hedge, trade} (idb);
\draw[arr,<->] (pd) -- node[above=11pt,font=\footnotesize,align=center] {quotes on request} (cli);
\draw[arr,<->] (fut) -- node[below=11pt,font=\footnotesize] {basis} (idb);
\draw[arr,dashed] (idb) -- node[below=11pt,font=\footnotesize] {trades cleared} (ccp);
\end{tikzpicture}

\omcaption{The structure of the US Treasury market. The Treasury sells to
dealers and direct bidders at auction; clients buy mostly from dealers on
request; dealers and principal trading firms trade with each other in the
inter-dealer order books and in futures. The clearing mandate
(\cref{dat:m2:the-treasury-market:clearing}) will send more of the
secondary trading through a clearing house.}\label{fig:m2:the-treasury-market:structure}
\end{omfigure}

\section{Principal trading firms and the electronic market}

Principal trading firms (\emph{proprietary trading firms}, One Quant Book 1,
chapter 1) trade their own capital, mostly electronically, mostly at short
horizons, and mostly in the benchmark securities and the futures. On 15
October 2014 they accounted for more than half of the volume in the
inter-dealer cash market and the futures that the regulators analysed, and
their activity was concentrated: the ten most active did more than 90\% of
all principal-trading-firm trading in the cash market. They provided most of
the displayed depth, and they took it away when they chose. A market whose
liquidity is supplied by firms without an obligation to supply it, the
dealers' auction obligation having no counterpart in the secondary market, is
deep in normal times and can be thin on the days it is needed.

\section{15 October 2014 and March 2020}

\begin{definition}[Flash rally]\label{def:m2:the-treasury-market:flash}
A \emph{flash rally}\index{flash rally} is a large, fast rise in prices (fall
in yields) that reverses within minutes without news to explain either move,
the mirror image of a flash crash (One Quant Book 1, chapter 31).
\end{definition}

The ten-year note's twelve minutes on 15 October 2014
(\cref{fig:m2:the-treasury-market:oct2014}) followed a weak retail sales
release at 8:30 that had already pushed yields down and depth out of the
books. In the event window the order books thinned to a fifth of their usual
depth in the ten-year, trading volume surged, and self-trading by single firms
was unusually high. The report found no single cause. The yield closed at
2.14\%, six basis points below the previous close, after a range of 37 basis
points, a range seen only three times since 1998 and each time after
significant news.

\begin{omfigure}
\begin{tikzpicture}[font=\small,x=0.9cm,y=0.9cm]
\draw[->,thick] (0,-0.9) -- (12.5,-0.9) node[right,font=\footnotesize] {time (ET)};
\draw[->,thick] (0,-0.9) -- (0,4.7) node[above,font=\footnotesize] {ten-year yield (\%)};
\foreach \y/\l in {0.5/1.85,1.75/1.95,3.0/2.05,4.25/2.15}{\draw (0.1,\y) -- (-0.1,\y) node[left,font=\footnotesize] {\l};}
% y = 0.5 + 12.5 (yield - 1.85); x = 1 + 0.85 (minutes after 9:33:19), so time is to scale
\coordinate (a) at (1.0,2.625);    % 2.02 at 9:33:19
\coordinate (b) at (6.383,0.625);  % 1.86 at 9:39:39 (6 min 20 s later)
\coordinate (c) at (10.578,2.25);  % 1.99 at 9:44:35 (4 min 56 s after that)
\draw[omDef,very thick] (a) -- (b) -- (c);
\foreach \p/\t/\v/\pos in {a/9:33:19/2.02/above,b/9:39:39/1.86/below,c/9:44:35/1.99/above}{
  \fill[omDef] (\p) circle (2.2pt);
  \node[\pos=5pt,font=\footnotesize,align=center] at (\p) {\t\\\v\%};}
\draw[omProp,dashed] (0,4.125) -- (12.3,4.125) node[below left,font=\footnotesize,text=omProp] {close 2.14\%};
\end{tikzpicture}

\omcaption{The event window of 15 October 2014: the three yields and times
reported by the regulators' joint staff report, joined by straight lines (the
path between them was not smooth), and the day's close. Sixteen basis points
down in six minutes, and back, in the benchmark of the dollar yield
curve. Data: Joint Staff Report, 2015.}\label{fig:m2:the-treasury-market:oct2014}
\end{omfigure}

\begin{definition}[Dash for cash]\label{def:m2:the-treasury-market:dash}
A \emph{dash for cash}\index{dash for cash} is a rush by investors of all
kinds to sell liquid assets, including the safest government bonds, to raise
cash, so that the securities normally bought in a crisis are sold instead.
\end{definition}

In March 2020, as the pandemic closed economies, investors of every kind
sought cash. The international review of the episode found that substantial
sales of Treasuries by some leveraged non-bank investors and by foreign
holders met dealers who could not, or would not, absorb them; the combination
became self-reinforcing, and prices that normally move together, \omterm{def:m2:the-treasury-market:otr}{on-the-run}
and \omterm{def:m2:the-treasury-market:otr}{off-the-run} notes, Treasuries and their futures, came apart. On 15 March the Federal Reserve cut its rate to zero
and announced purchases of at least USD~500 billion of Treasuries; on 23 March
it changed the instruction to purchases ``in the amounts needed'' for the
market to function. Its holdings rose from USD~2\,523 billion on 11 March to
USD~3\,789 billion on 15 April, 1\,266 billion in five weeks
(\cref{fig:m2:the-treasury-market:fed2020}). The central bank became, for a
month, the market's dealer of last resort, the role the clearing mandate and
the regulators' other reforms since then are meant to make less necessary.

\begin{omfigure}
\begin{tikzpicture}
\begin{axis}[width=11cm,height=4.8cm,
  ylabel={USD billion},
  tick label style={font=\small},label style={font=\small},
  xmin=0,xmax=238,ymin=2000,ymax=4500,grid=major,grid style={gray!25},
  xtick={28,59,88,119,149,180,210},xticklabels={Jan,Feb,Mar,Apr,May,Jun,Jul},
  table/col sep=comma]
\addplot[omDef,very thick,mark=*,mark size=1.1pt] table[x=t,y=usdbn]{figdata/markets-2/04-the-treasury-market/fed2020.csv};
\draw[omThm,dashed] (axis cs:102,2000) -- (axis cs:102,4500);
\draw[omThm,dashed] (axis cs:110,2000) -- (axis cs:110,4500);
\node[font=\footnotesize,anchor=east,text=omThm] at (axis cs:101,4250) {15 Mar};
\node[font=\footnotesize,anchor=west,text=omThm] at (axis cs:111,4250) {23 Mar};
\end{axis}
\end{tikzpicture}

\omcaption{Treasury securities held by the Federal Reserve, weekly, December
2019 to July 2020. The two dashed lines mark the announcement of at least USD~500
billion of purchases and the change to purchases in the amounts needed. Data:
Federal Reserve H.4.1, via FRED series TREAST.}\label{fig:m2:the-treasury-market:fed2020}
\end{omfigure}

\begin{dated}{2026-09}{The US Treasury clearing mandate}\label{dat:m2:the-treasury-market:clearing}
Adopted by the SEC in December 2023, the rule requires members of Treasury
clearing agencies to clear eligible secondary-market trades: eligible cash
transactions from 31 December 2026 and eligible repo from 30 June 2027 (dates
extended by a year in February 2025). On 22 September 2026 an SEC commissioner
said the SEC did not intend to extend them again. Two clearing agencies were
approved to act as central counterparties for Treasury cash and repo trades.
\end{dated}

\section{Tutorial: running an auction}\label{tut:m2:the-treasury-market:auction}

\begin{tutorial}
\textbf{Goal.} Run the synthetic ten-year auction of
\cref{ex:m2:the-treasury-market:auction} from its bid book and compute every
statistic the market reads at 13:01. \textbf{End state:}
\cref{fig:m2:the-treasury-market:demand} and the numbers of the example.
\begin{tutsteps}
\item \textbf{Cap and sort.} Each bidder's bids count up to 35\% of the offering.
\omcode{code/firm/tsyauction/firm_tsyauction.py}{33}{42}{Recognising each bidder's bids up to the cap.}\label{lst:m2:the-treasury-market:cap}
\item \textbf{Fill to the stop-out}, with the pro-rata percentage rounded up as
the regulation says.
\omcode{code/firm/tsyauction/firm_tsyauction.py}{45}{65}{The uniform-price auction: stop-out yield, allotment at the stop, awards.}\label{lst:m2:the-treasury-market:run}
\item \textbf{Read the result.} \texttt{tsy\_auction\_demo.summary()} gives the
stop-out 4.198\%, the tail of 1.8 basis points against the when-issued 4.180\%,
the \omterm{def:m2:the-treasury-market:tail}{bid-to-cover ratio} 1.82, the allotment at the stop 95.02\% and the award
shares; \texttt{coupon\_from\_high\_yield(0.04198)} gives 4.125.
\item \textbf{Draw the demand curve} with \texttt{fig\_tsy\_auction.py}.
\end{tutsteps}
\textbf{What to change next.} Move every indirect bid one basis point higher
and see how far the stop-out moves (\cref{exo:m2:the-treasury-market:7}); then
give one bidder 60\% of the book and watch the cap bite.
\end{tutorial}

\section{Build: the auction analyser}\label{bld:m2:the-treasury-market:auction}

\begin{build}
\bfield{Purpose} The miniature firm's rates desk bids in auctions and trades
around them; its research reads every auction's statistics. This component
runs an auction from a bid book and scores a result against the when-issued
market.
\bfield{Interface} \texttt{Bid(bidder, yld, amount)};
\texttt{run\_auction(offering, noncompetitive, bids)} returning
\texttt{Result(stop, allotment\_at\_stop, awards, competitive\_tendered,
bid\_to\_cover)}; \texttt{cap\_bids}; \texttt{tail\_bp(stop, when\_issued)};
\texttt{cumulative\_demand}; \texttt{coupon\_from\_high\_yield}.
\bfield{Rules} Noncompetitive first; 35\% cap per bidder; stop-out at the
yield where cumulative bids reach the remainder; pro-rata percentage rounded
up to 0.01\%; an uncovered auction is an error; coupon rounded down to an
eighth, minimum one eighth.
\bfield{Acceptance tests} \texttt{code/firm/tsyauction/tests/}: every winner
pays the stop-out; noncompetitive awards reduce the competitive capacity; the
cap; an uncovered book; the sign of the tail; the coupon rule.
\bfield{Stretch} Read the Treasury's published auction results and compute
the tail against a when-issued feed; add the bidder categories the Treasury
publishes; price the new note with the bond library
(\cref{bld:m2:government-bonds:bond}).
\end{build}

\begin{omsources}
\item US Treasury, \emph{How Auctions Work} and \emph{Auctions In Depth}
(TreasuryDirect); 31 CFR 356.20.
\item Federal Reserve Bank of New York, \emph{Primary Dealers}.
\item US Department of the Treasury, Board of Governors of the Federal
Reserve System, Federal Reserve Bank of New York, SEC and CFTC, \emph{Joint
Staff Report: The U.S. Treasury Market on October 15, 2014}, July 2015.
\item Financial Stability Board, \emph{Holistic Review of the March Market
Turmoil}, 17 November 2020.
\item Federal Reserve press release, 15 March 2020; Federal Reserve Bank of
New York, \emph{FAQs: Treasury Purchases}, 23 March 2020; H.4.1 via FRED.
\item SEC, press release 2025-43; M.\,T.\ Uyeda, remarks at the 2026 U.S.
Treasury Market Conference, 22 September 2026.
\end{omsources}

\section{Exercises}

\begin{exercise}[$\star$]\label{exo:m2:the-treasury-market:1}
USD~30 billion are offered; noncompetitive bids total 3 billion. Four bidders
bid 10 billion at 4.100\%, 8 billion at 4.105\%, 10 billion at 4.110\% and 6
billion at 4.115\%. Give the stop-out yield, the allotment at the stop, and
the \omterm{def:m2:the-treasury-market:tail}{bid-to-cover ratio}.
\end{exercise}

\begin{exercise}[$\star$]\label{exo:m2:the-treasury-market:2}
The when-issued yield at 13:00 is 4.285\%. Give the tail if the stop-out is
4.292\%, and if it is 4.281\%. Which result would the market call strong?
\end{exercise}

\begin{exercise}[$\star$]\label{exo:m2:the-treasury-market:3}
Give the coupon of a new note whose auction stops at 4.198\%, and at 3.999\%.
Is the new note issued above or below par?
\end{exercise}

\begin{exercise}[$\star\star$]\label{exo:m2:the-treasury-market:4}
An \omterm{def:m2:the-treasury-market:otr}{on-the-run} ten-year yields 4.18\% and the previous ten-year, three months
older, 4.21\%. Give three reasons for the gap and describe the trade that
bets on it closing. What can go wrong?
\end{exercise}

\begin{exercise}[$\star\star$]\label{exo:m2:the-treasury-market:5}
In \cref{ex:m2:the-treasury-market:auction}, what is a \omterm{def:m2:the-treasury-market:pd}{primary dealer}'s
pro-rata share of the offering? The dealers took 38.3\% of the auction. What
does that number tell the market that the \omterm{def:m2:the-treasury-market:tail}{bid-to-cover ratio} does not?
\end{exercise}

\begin{exercise}[$\star\star$]\label{exo:m2:the-treasury-market:6}
A single investor wants USD~20 billion of a 42-billion auction. How much of its
bid is recognised? How could it still end up owning 20 billion?
\end{exercise}

\begin{exercise}[$\star\star\star$]\label{exo:m2:the-treasury-market:7}
\emph{Coding.} In \texttt{tsy\_auction\_demo.bid\_book()}, move every indirect
bidder's yield one basis point higher. Report the new stop-out and tail. Why
does the stop-out move by less than, or as much as, the shift?
\end{exercise}

\begin{exercise}[$\star\star\star$]\label{exo:m2:the-treasury-market:8}
\emph{Find the flaw.} ``In a \omterm{def:m2:the-treasury-market:auction}{uniform-price auction}, bid a yield a little above
what you think the note is worth, because you pay what you bid.'' Correct the
statement. Who does have a reason to shade their bids, and why?
\end{exercise}

\section{Problem: Tailing}

\begin{problem}[{Weekend problem --- a primary dealer's auction}]\label{pb:m2:the-treasury-market:1}
\omterm{def:m2:the-treasury-market:pd}{Primary dealer} D07 takes part in the auction of \cref{ex:m2:the-treasury-market:auction}.
At 12:59 it has sold USD~1 billion of the new note when-issued at 4.180\% to
clients, and it bids USD~500 million at 4.188\%, 500 million at 4.195\% and 615
million at 4.205\%. Its round-trip costs are a quarter of a basis point of
DV01. The new note's DV01 at the stop-out is USD~807 per million.

\medskip\noindent\textbf{Part I --- The auction.}
\begin{enumerate}
\item How much is available to competitive bidders?
\item Give the stop-out yield and the tail.
\item Give the \omterm{def:m2:the-treasury-market:tail}{bid-to-cover ratio} and the allotment at the stop.
\item Give the shares of indirect bidders, direct bidders and dealers.
\item Give the coupon and the price every winner pays.
\end{enumerate}

\medskip\noindent\textbf{Part II --- The dealer.}
\begin{enumerate}[resume]
\item What is D07's pro-rata share of the offering? Does its bid meet the
obligation?
\item Which of its bids are filled, and how much is it awarded?
\item What is the DV01 of its award?
\item It is short 1 billion when-issued at 4.180\% and long 1 billion at the
stop-out. What is its P\&L after costs?
\item What would it have been had the auction stopped through by 0.2 basis
points (at 4.178\%), with the same award?
\end{enumerate}

\medskip\noindent\textbf{Part III --- Reading the result.}
\begin{enumerate}[resume]
\item At what tail does the dealer's position break even?
\item Why does the market read a large tail as weak demand, and what does the
when-issued note do in the minute after the result?
\item Why did the dealer bid part of its share above the when-issued yield
at all?
\item Had the indirect bidders all bid one basis point higher, where would
the auction have stopped?
\item What would the dealer have wished to do differently, knowing the tail
in advance, and why can it not?
\end{enumerate}

\medskip\noindent\textbf{Part IV --- Judgement.}
\begin{enumerate}[resume]
\item Why does the Treasury use a uniform-price rather than a
discriminatory (pay-your-bid) auction?
\item Why do the clients who bought when-issued from D07 not simply bid in
the auction themselves?
\item What does the dealer's auction obligation cost it in a weak auction,
and what does it get in return?
\item State the \emph{named result}: the dealer's P\&L per basis point of tail
on its award, its P\&L in this auction, and the break-even tail.
\item In one sentence: what is a tail?
\end{enumerate}
\end{problem}

\section{Interview questions}

\begin{interviewq}[$\star$ \iqroles{trader, researcher}]\label{iq:m2:the-treasury-market:1}
What is an \omterm{def:m2:the-treasury-market:tail}{auction tail}, and what does a large one tell you?
\end{interviewq}

\begin{interviewq}[$\star$ \iqroles{trader, bank}]\label{iq:m2:the-treasury-market:2}
Why does an \omterm{def:m2:the-treasury-market:otr}{on-the-run} Treasury usually yield less than an \omterm{def:m2:the-treasury-market:otr}{off-the-run} of
almost the same maturity?
\end{interviewq}

\begin{interviewq}[$\star\star$ \iqroles{trader, researcher}]\label{iq:m2:the-treasury-market:3}
How does a \omterm{def:m2:the-treasury-market:auction}{uniform-price auction} change a bidder's strategy compared with a
pay-your-bid auction?
\end{interviewq}

\begin{interviewq}[$\star\star$ \iqroles{researcher, developer}]\label{iq:m2:the-treasury-market:4}
What happened in the Treasury market on 15 October 2014, and what would you
look for in order-book data to understand a day like it?
\end{interviewq}

\begin{interviewq}[$\star\star$ \iqroles{trader, bank}]\label{iq:m2:the-treasury-market:5}
In March 2020 Treasuries sold off while equities crashed. Why did the safest
asset fall?
\end{interviewq}

\begin{interviewq}[$\star\star\star$ \iqroles{developer, researcher}]\label{iq:m2:the-treasury-market:6}
Design the data model and the tests for a system that ingests Treasury
auction results and computes each auction's tail in real time.
\end{interviewq}
